%! Polynomial Functions and Complex Zeros — Unit 2, Lesson 2.4 — Homework
\documentclass[10pt]{article}
\usepackage{precalculus-article}
\usepackage{precalculus-boxes}
\newcommand{\TallMath}[1]{$\displaystyle #1\rule[-1.4em]{0pt}{3.2em}$}

\begin{document}

\pageheader{Unit 2, Lesson 2.4}{Homework}
\namedateperiod

\vspace{0.1in}

\begin{enumerate}[leftmargin=1.5em, itemsep=12pt]

\item Write a polynomial with real coefficients of \emph{minimum degree} that has zeros
      at $x = 3$ (multiplicity 2) and $x = 1 - 2i$. Leave your answer in factored form.
      State the degree.

      \vspace{0.06in}
      $p(x) = $ \blank{8cm} \hfill Degree: \blank{2cm}

\item A degree-5 polynomial with real coefficients has zeros at $x = -2$,
      $x = 0$ (multiplicity 2), and $x = 1 + i$.
      List \emph{all five} zeros.

      \writeline

\item Let $p(x) = x^4 - 6x^2 + 9$.
      \begin{enumerate}[leftmargin=1.5em, itemsep=6pt]
        \item Factor $p(x)$ completely. (\textit{Hint: is it a perfect square?})
              \writelines{2}
        \item List all zeros and their multiplicities.
              \writeline
        \item At each real zero, does the graph cross or is it tangent to the $x$-axis?
              \writeline
      \end{enumerate}

\item Is $g(x) = 2x^3 - 5x$ even, odd, or neither? Show the complete algebraic test.

      \vspace{0.06in}
      $g(-x) = $ \blank{7cm}

      \vspace{0.08in}
      Simplified: \blank{7cm}

      \vspace{0.08in}
      $g$ is (circle): \textbf{Even} \quad \textbf{Odd} \quad \textbf{Neither}

\item \textbf{AP-Style Multiple Choice.} A polynomial $p(x)$ is defined by
      \[
        p(x) = (x+3)^2(x-1)(x-5)^3.
      \]
      Which of the following correctly describes all zeros of $p$ and the graph's behavior
      at each?
      \begin{enumerate}[leftmargin=1.5em, itemsep=3pt, label=(\Alph*)]
        \item $x=-3$ (tangent), $x=1$ (crosses), $x=5$ (crosses)
        \item $x=-3$ (crosses), $x=1$ (tangent), $x=5$ (crosses)
        \item $x=-3$ (tangent), $x=1$ (crosses), $x=5$ (tangent)
        \item $x=3$ (tangent), $x=-1$ (crosses), $x=-5$ (crosses)
      \end{enumerate}

\end{enumerate}

\vspace{0.1in}

\begin{extensionbox}
A polynomial $p$ of unknown degree has output values recorded at equally-spaced inputs.
The third differences of those outputs are all constant, each equal to $6$.
\begin{enumerate}[leftmargin=1.5em, itemsep=5pt]
  \item What is the degree of $p$? Explain how successive differences reveal degree.
        \writeline
  \item Write a general form of a polynomial of that degree: \blank{5cm}
\end{enumerate}
\end{extensionbox}

\vspace{0.1in}

\begin{reflectionbox}
\textbf{Preview:} In Lesson 1.6 we will study rational functions. As you finish this
homework, think about: if a polynomial can be written as a ratio $\dfrac{p(x)}{q(x)}$,
what happens when $q(x) = 0$? We'll explore vertical asymptotes and holes next class.
\end{reflectionbox}

\end{document}
